\documentclass[12pt,a4paper]{article}
% Drake_preamble.tex - LaTeX Preamble Template
% 作者: 謝慕揚, MD, PhD, FESC
% 
% 使用說明:
% 1. 表格請使用 longtable，不要有垂直分隔線 |
% 2. 表格內換行使用 \newline，不要使用 \\
% 3. 藥物名稱、檢驗、檢查請維持英文
% 4. Reference 格式請使用 \href 加上驗證過的 DOI

% Including essential packages for Chinese typesetting
\usepackage{xeCJK}
\usepackage{fontspec}
% Configuring Chinese font with Noto Serif CJK TC, fallback to SimSun
\IfFontExistsTF{Noto Serif CJK TC}{
  \setCJKmainfont{Noto Serif CJK TC}
  \setCJKsansfont{Noto Serif CJK TC}
  \setCJKmonofont{Noto Serif CJK TC}
}{\setCJKmainfont{SimSun}
  \setCJKsansfont{SimSun}
  \setCJKmonofont{SimSun}
}

% Including other necessary packages
\usepackage{geometry}
\usepackage{titlesec}
\usepackage{enumitem}
\usepackage{xcolor}
\usepackage{colortbl}
\usepackage{tcolorbox}
\usepackage{booktabs}
\usepackage{longtable}
\usepackage{array}
\usepackage{graphicx}
\usepackage{tikz}
\usepackage{fancyhdr}
\usepackage{multicol}
\usepackage{datetime2}
\usepackage{hyperref}
\usepackage{mdframed}
\usepackage{amsmath}
\usepackage{amssymb}
\usepackage{multirow}

\hypersetup{
    colorlinks=true,
    linkcolor=emergency_blue,
    citecolor=emergency_blue,
    urlcolor=emergency_blue,
    pdfborder={0 0 0}
}

% Defining page geometry
\geometry{
  top=2.5cm,
  bottom=2.5cm,
  left=2cm,
  right=2cm,
  headheight=25pt
}

% Adjusting topmargin to compensate for increased headheight
\addtolength{\topmargin}{-10pt}

% Defining colors
\definecolor{emergency_red}{RGB}{186,24,27}
\definecolor{emergency_blue}{RGB}{0,114,188}
\definecolor{emergency_gray}{RGB}{120,120,120}
\definecolor{highlight_yellow}{RGB}{255,250,205}

% Setting title formats
\titleformat{\section}[block]{\Large\bfseries\color{emergency_red}}{\thesection}{10pt}{}
\titleformat{\subsection}[block]{\large\bfseries\color{emergency_blue}}{\thesubsection}{8pt}{}
\titleformat{\subsubsection}[block]{\normalsize\bfseries\color{emergency_gray}}{\thesubsubsection}{6pt}{}

% Defining custom environment for important notes
\newmdenv[
    linecolor=emergency_red,
    backgroundcolor=highlight_yellow,
    linewidth=2pt,
    topline=true,
    bottomline=true,
    leftline=true,
    rightline=true,
    innertopmargin=10pt,
    innerbottommargin=10pt,
    innerleftmargin=10pt,
    innerrightmargin=10pt
]{importantbox}

% Defining note command with tcolorbox
\newcommand{\note}[1]{
    \par
    \noindent
    \begin{tcolorbox}[
        colback=emergency_blue!5!white,
        colframe=emergency_blue,
        title=\textbf{重點提醒},
        fonttitle=\bfseries,
        boxrule=0.5pt,
        arc=3pt,
        left=6pt,
        right=6pt,
        top=6pt,
        bottom=6pt
    ]
    #1
    \end{tcolorbox}
}

% Key findings box
\newcommand{\keyfinding}[1]{
    \par
    \noindent
    \begin{tcolorbox}[
        colback=yellow!10!white,
        colframe=orange!75!black,
        title=\textbf{核心發現摘要},
        fonttitle=\bfseries,
        boxrule=0.5pt,
        arc=3pt
    ]
    #1
    \end{tcolorbox}
}

% Defining custom date format for ISO 8601
\DTMnewdatestyle{iso}{
  \renewcommand{\DTMdisplaydate}[4]{\number##1-\number##2-\number##3}
  \renewcommand{\DTMDisplaydate}[4]{\number##1-\number##2-\number##3}
}
\DTMsetdatestyle{iso}
\newcommand{\mydate}{\DTMtoday}


% Custom header for this document
\pagestyle{fancy}
\fancyhf{}
\fancyhead[L]{\textcolor{emergency_red}{Intensive Care Medicine 教學講義}}
\fancyhead[R]{\textcolor{emergency_blue}{Broad-Spectrum Antibiotics}}
\fancyfoot[C]{\textcolor{emergency_gray}{\thepage}}
\renewcommand{\headrulewidth}{1pt}
\renewcommand{\headrule}{\hbox to\headwidth{\color{emergency_red}\leaders\hrule height \headrulewidth\hfill}}

\begin{document}

\begin{center}
{\LARGE\bfseries\textcolor{emergency_red}{Intensive Care Medicine Editorial}}\\[0.5cm]
{\Large\textbf{When Should I Start Broad-Spectrum Antibiotics?}}\\[0.3cm]
{\large 何時應該開始使用廣效抗生素？}\\[0.5cm]
{\normalsize Intensive Care Med 2024;50:1908-1911. DOI: 10.1007/s00134-024-07654-7}\\[0.3cm]
{\normalsize 整理：謝慕揚 MD, PhD, FESC \quad 日期：\mydate}
\end{center}

\keyfinding{
\textbf{核心訊息：}Early targeted therapy 基於培養及/或分子檢測結果，可能是最佳的廣效抗生素使用策略。不必要的廣效抗生素與較差預後相關。

\textbf{臨床建議：}
\begin{itemize}
    \item \textbf{MDRO 風險評估：}根據地理、臨床、微生物風險因子分層
    \item \textbf{快速診斷工具：}利用分子檢測輔助決策，注意 NPV 及 PPV
    \item \textbf{De-escalation：}臨床改善及培養結果出來後，應在 24-72 小時內進行抗生素合理化及降階
\end{itemize}
}

\tableofcontents
\newpage

%==============================================================
\section{文獻概述}
%==============================================================

本篇為 Intensive Care Medicine 的 Editorial，由法國 Bichat Hospital 的 Timsit 等人撰寫，討論重症病人廣效抗生素使用時機的決策考量。

\subsection{作者資訊}
\textbf{通訊作者：}Jean-François Timsit, MD\\
\textbf{機構：}Medical and Infectious Diseases ICU (MI2), AP-HP, Bichat Hospital, Paris, France; Université Paris-Cité, IAME, INSERM, Paris, France

%==============================================================
\section{背景}
%==============================================================

\subsection{廣效抗生素的定位}

\begin{itemize}
    \item 常被視為治療嚴重或危及生命感染的「magic bullets」
    \item 目的：在培養結果出來前，涵蓋 Gram-positive、Gram-negative 及 resistant pathogens
    \item 在缺乏快速診斷工具的環境中被視為救命藥物
\end{itemize}

\subsection{ICU 中的使用問題}

\begin{itemize}
    \item 緊急治療重症病人的需求常導致 \textbf{overuse}
    \item Inadequate therapy 的延遲與較高 morbidity 及 mortality 相關
    \item 但不加選擇地使用廣效抗生素顯著促進了 \textbf{antimicrobial resistance (AMR)} 的發展
\end{itemize}

%==============================================================
\section{廣效抗生素的潛在益處}
%==============================================================

\note{
\textbf{使用廣效抗生素的適當情境：}
\begin{itemize}
    \item Sepsis 及嚴重感染（如 ventilated hospital-acquired pneumonia, vHAP）
    \item 高 MDRO 風險且有 shock（容錯空間較小）
    \item 快速 de-escalation 配合培養結果可減少 antibiotic selection pressure
\end{itemize}
}

\subsection{需使用廣效抗生素進行 Targeted Therapy 的情況}

\begin{longtable}{p{5cm}p{9cm}}
\toprule
\textbf{感染/病原體} & \textbf{建議抗生素} \\
\midrule
\endfirsthead

\toprule
\textbf{感染/病原體} & \textbf{建議抗生素} \\
\midrule
\endhead

\bottomrule
\endfoot

3rd-generation cephalosporin-resistant Enterobacterales bacteremia & Meropenem \\
\midrule
Carbapenem-resistant Enterobacterales severe infections & Ceftazidime-avibactam\newline Meropenem-vaborbactam \\
\midrule
Carbapenem-resistant \textit{Pseudomonas aeruginosa} vHAP & Ceftolozane-tazobactam \\
\midrule
Community-acquired CNS infections & Ceftriaxone（優於 amoxicillin）\\
\midrule
Nosocomial CNS infections & Meropenem（優於 piperacillin-tazobactam）\\
\bottomrule
\end{longtable}

%==============================================================
\section{廣效抗生素的負面影響}
%==============================================================

\subsection{AMR 以外的危害}

\begin{itemize}
    \item 增加照護成本
    \item 藥物毒性
    \item \textit{Clostridioides difficile} colitis
    \item Acute kidney injury
\end{itemize}

\subsection{Webb et al. 研究：Community-Acquired Pneumonia}

\begin{longtable}{p{5cm}p{9cm}}
\toprule
\textbf{項目} & \textbf{發現} \\
\midrule
\endfirsthead

\toprule
\textbf{項目} & \textbf{發現} \\
\midrule
\endhead

\bottomrule
\endfoot

\textbf{研究對象} & 1,995 位 community-acquired pneumonia 病人 \\
\midrule
\textbf{廣效抗生素使用率} & 39.7\% \\
\midrule
\textbf{MDRO 感染率} & 僅 3\% \\
\midrule
\textbf{結果} & 廣效抗生素與 \textbf{increased mortality} 相關\newline
17.5\% 死亡者有 antibiotic-associated adverse events \\
\bottomrule
\end{longtable}

\subsection{Rhee et al. 研究：Community-Acquired Sepsis}

\begin{longtable}{p{5cm}p{9cm}}
\toprule
\textbf{項目} & \textbf{發現} \\
\midrule
\endfirsthead

\toprule
\textbf{項目} & \textbf{發現} \\
\midrule
\endhead

\bottomrule
\endfoot

\textbf{研究對象} & 15,183 位有細菌學記錄的 community-acquired sepsis 成人 \\
\midrule
\textbf{Appropriate antibiotics} & 81.6\% \\
\midrule
\textbf{MDRO 感染比例} & <30\% \\
\midrule
\textbf{不必要的廣效抗生素} & 67.8\%（涵蓋 MRSA、VRE、ceftriaxone-resistant GNB，但均未分離出）\\
\midrule
\textbf{Hospital death OR} & \textbf{1.27 (1.06-1.4)} 相較於 appropriate broad-spectrum \\
\bottomrule
\end{longtable}

%==============================================================
\section{MDRO 風險因子評估}
%==============================================================

\subsection{風險分層}

\begin{longtable}{p{4cm}p{10cm}}
\toprule
\textbf{風險類型} & \textbf{風險因子} \\
\midrule
\endfirsthead

\toprule
\textbf{風險類型} & \textbf{風險因子} \\
\midrule
\endhead

\bottomrule
\endfoot

\textbf{地理風險} &
• MDR endemic 地區（+）\newline
• 近期至 MDR endemic region 旅行（+）\newline
• 近期在 MDR endemic region 醫療暴露（++）\\
\midrule
\textbf{臨床風險} &
• 過去 6 個月住院 >2 天（+）\newline
• 目前住院 >5 天（++）\newline
• 抗生素暴露（+）\newline
• >1 種抗生素暴露（++）\newline
• Invasive devices（++）\newline
• Dependency ADL（+）\\
\midrule
\textbf{微生物風險} & 曾被 MDR pathogen colonization/infection（++）\\
\midrule
\textbf{不良預後風險} &
• Septic shock, ARDS（+）\newline
• Neutropenia（+）\newline
• Other immunosuppression（+）\\
\bottomrule
\end{longtable}

\subsection{風險分層與抗生素選擇}

\note{
\textbf{決策原則：}
\begin{itemize}
    \item \textbf{0 MDR risk factors → Low MDR risk：}Narrow spectrum，針對 core pathogens
    \item \textbf{1 MDR risk factor → Intermediate MDR risk：}考慮 MDR pathogen coverage，利用 early diagnostics NPV
    \item \textbf{≥2 MDR risk factors → High MDR risk：}Broad spectrum，針對 core pathogens 及 MDR pathogens，利用 NPV 及 PPV
\end{itemize}
}

%==============================================================
\section{快速診斷工具的角色}
%==============================================================

\subsection{應用時機}

\begin{itemize}
    \item 風險因子評估後（preanalytical phase）
    \item 在適當檢體上進行快速診斷可以：
    \begin{itemize}
        \item 精準化 targeted therapy
        \item 避免不必要的廣效抗生素
    \end{itemize}
    \item 可在數小時內獲得結果：urinary、respiratory、digestive、neurologic infections
\end{itemize}

\subsection{Pretest Probability 與檢測結果解讀}

\begin{longtable}{p{4cm}p{5cm}p{5cm}}
\toprule
\textbf{Pretest Probability} & \textbf{陰性結果} & \textbf{陽性結果} \\
\midrule
\endfirsthead

\toprule
\textbf{Pretest Probability} & \textbf{陰性結果} & \textbf{陽性結果} \\
\midrule
\endhead

\bottomrule
\endfoot

\textbf{Low} & Very high NPV → 可排除 MDRO & PPV 可能不足以接受為真陽性 \\
\midrule
\textbf{High} & NPV 可能不足以排除 MDRO & 應涵蓋 MDRO \\
\midrule
\textbf{有 Shock 或快速惡化} & NPV 可能不足以使用 narrow spectrum & 應涵蓋 MDRO \\
\bottomrule
\end{longtable}

\subsection{Multiplex PCR 的限制}

\begin{itemize}
    \item 範例：在 ICU 嚴重肺炎中，若 OXA-48 Enterobacterales 風險為 0.2\%
    \item Multiplex PCR 不會漏診，但會在 \textbf{37.5\% 的案例中過度診斷}
    \item Syndromic mPCR panels 的多重目標應謹慎解讀 pretest probabilities
\end{itemize}

%==============================================================
\section{抗生素 De-Escalation}
%==============================================================

\note{
\textbf{De-Escalation 時機：}
\begin{itemize}
    \item 臨床狀態改善
    \item 治療前採集的培養結果出來
    \item 應在 \textbf{24-72 小時內}進行 antimicrobial rationalization 及 de-escalation
\end{itemize}
}

%==============================================================
\section{決策考量因素}
%==============================================================

\begin{longtable}{p{5cm}p{9cm}}
\toprule
\textbf{考量因素} & \textbf{說明} \\
\midrule
\endfirsthead

\toprule
\textbf{考量因素} & \textbf{說明} \\
\midrule
\endhead

\bottomrule
\endfoot

\textbf{Guidelines} & 熟悉國家及機構指引 \\
\midrule
\textbf{Infection Source} & 根據感染來源識別常見病原體 \\
\midrule
\textbf{Prior MDRO Carriage} & 過去的 MDRO colonization 或感染史 \\
\midrule
\textbf{Prior Antimicrobial Exposure} & 先前的抗生素使用史 \\
\midrule
\textbf{Local Epidemiology} & 了解國家、醫院、ICU 的 MDRO epidemiology 及 resistance patterns \\
\midrule
\textbf{Clinical Severity} & 根據病人臨床嚴重程度個人化決策 \\
\midrule
\textbf{Cost-Benefit Ratio} & 評估每種策略的預估 cost-benefit ratio \\
\bottomrule
\end{longtable}

%==============================================================
\section{Take-Home Message}
%==============================================================

\begin{enumerate}
    \item \textcolor{red}{\textbf{最佳策略：}}Early targeted therapy 基於培養及/或分子檢測結果，可能是最佳的廣效抗生素使用策略
    \item \textcolor{red}{\textbf{風險評估：}}徹底了解 guidelines、根據感染來源識別常見病原體、評估 MDRO carriage 及 antimicrobial exposure 史、了解 local epidemiology
    \item \textcolor{red}{\textbf{快速診斷：}}利用微生物實驗室及快速診斷工具，可減少經驗性廣效抗生素使用或允許快速 de-escalation
    \item \textcolor{red}{\textbf{Stewardship Team：}}建議由 multidisciplinary antibiotic stewardship team 協助此過程
\end{enumerate}

%==============================================================
\section{縮寫對照表}
%==============================================================

\begin{longtable}{p{3cm}p{11cm}}
\toprule
\textbf{縮寫} & \textbf{全名} \\
\midrule
\endfirsthead

\toprule
\textbf{縮寫} & \textbf{全名} \\
\midrule
\endhead

\bottomrule
\endfoot

AMR & Antimicrobial Resistance (抗微生物藥物抗藥性) \\
MDRO & Multidrug-Resistant Organism (多重抗藥性病原體) \\
vHAP & Ventilated Hospital-Acquired Pneumonia \\
MRSA & Methicillin-Resistant \textit{Staphylococcus aureus} \\
VRE & Vancomycin-Resistant Enterococci \\
GNB & Gram-Negative Bacteria \\
CNS & Central Nervous System \\
NPV & Negative Predictive Value \\
PPV & Positive Predictive Value \\
mPCR & Multiplex Polymerase Chain Reaction \\
ADL & Activities of Daily Living \\
ARDS & Acute Respiratory Distress Syndrome \\
ESBL & Extended-Spectrum Beta-Lactamase \\
\bottomrule
\end{longtable}

%==============================================================
\section{參考文獻}
%==============================================================

\begin{enumerate}
    \item Timsit JF, Depuydt P, Kanj SS. When should I start broad-spectrum antibiotics? \href{https://doi.org/10.1007/s00134-024-07654-7}{\textit{Intensive Care Med}. 2024;50:1908-1911.}

    \item Rhee C, Kadri SS, Dekker JP, et al. Prevalence of antibiotic-resistant pathogens in culture-proven sepsis and outcomes associated with inadequate and broad-spectrum empiric antibiotic use. \href{https://doi.org/10.1001/jamanetworkopen.2020.2899}{\textit{JAMA Netw Open}. 2020;3:e202899.}

    \item Webb BJ, Sorensen J, Jephson A, et al. Broad-spectrum antibiotic use and poor outcomes in community-onset pneumonia: a cohort study. \href{https://doi.org/10.1183/13993003.00057-2019}{\textit{Eur Respir J}. 2019;54:1900057.}

    \item Timsit JF, Bassetti M, Cremer O, et al. Rationalizing antimicrobial therapy in the ICU: a narrative review. \href{https://doi.org/10.1007/s00134-018-5382-9}{\textit{Intensive Care Med}. 2019;45:172-189.}

    \item Zakhour J, Haddad SF, Kerbage A, et al. Diagnostic stewardship in infectious diseases: a continuum of antimicrobial stewardship in the fight against antimicrobial resistance. \href{https://doi.org/10.1016/j.ijantimicag.2023.106816}{\textit{Int J Antimicrob Agents}. 2023;62:106816.}
\end{enumerate}

\end{document}
